\section{Discussion and limitations}\label{sec:discussion}

\paragraph{What is predictable} The timing of a swing's end is predictable out of sample, and most
of what is predictable follows from how a swing is defined. A swing that has retraced most of the
way to its threshold is close to being confirmed, whatever happens in the economy. The distance to
the threshold carries \MechShareDIST\% of the gain over Kaplan--Meier, and a flexible function of
age and distance matches the probability accuracy of every richer linear model. We read this as a
warning for the literature on predicting turning points: an accuracy figure says little unless it
is compared with what the labelling rule implies by itself. The random-walk null makes the same
point from another side. Rolling a stock's own recent returns through the rule, in random order,
already gives C~=~\CRW.

\paragraph{What the features add} Beyond the mechanics, the stock's recent returns, volatility,
price-limit hits and the state of the local market improve the ranking of swings consistently,
in \YearsAboveABCxMFlex{} of 10 test years against the flexible baseline. They do not make the
predicted probabilities more accurate. A trader who only needs to know which open swings are the
most fragile gains a little; a risk manager who needs calibrated probabilities gains nothing that
age and distance do not already give. This difference between ranking and calibration is why we
required both metrics to improve before counting a gain.

\paragraph{Why outside information does not help} Overnight US returns, volatility, rates,
currencies, commodities and geopolitical risk add nothing to the timing of reversals, and neither
do bank fundamentals or news attention. The signs are often as expected: a fall in US technology
stocks overnight raises the hazard that a Vietnamese up-swing ends in every yearly fit. The
effects are small and, we think, already in the prices that our features use. Foreign information
is priced at the open \citep{HamaoMasulisNg1990}, so by the time a swing's state is measured, the
global news of the night is part of it. News reports a bank's results before its statement becomes usable
under our timing rule, and the news count itself is mostly attention that the price already
reflects. The null result for block E holds when the statements are dated with margins of 0 or 20
sessions, so it does not come from our conservative timing rule.

\paragraph{Why our result differs from earlier survival studies} \citet{Yao2005} found no out-of-sample skill in
runs of daily index changes. Their unit ends on the first opposite day, so most of its variation is
daily noise. A volatility-scaled swing ends only after a move that is large for the stock, and the
distance to that threshold is an informative state variable that a daily run does not have. The
gain we find is real, but it comes from the unit of analysis as much as from the models.

\paragraph{No trading value} The ranking gain does not survive the frictions of the Vietnamese
market. Long-only trading, settlement delays and price limits leave little room for an exit rule to
add value over simply following the swings, and the model exit trails buy-and-hold over a decade in
which the market rose. This is consistent with weak-form efficiency after costs.

\paragraph{Limitations} Several limits qualify these results.
\begin{itemize}
\item \emph{One market.} The sample covers VN30 stocks and listed banks in Vietnam. A market with
  different price limits, settlement or investor base could behave differently.
\item \emph{One labelling rule.} Swings follow one directional-change rule, with three threshold
  multiples. Other definitions of a trend, such as moving-average crossings, could give different
  answers, although the mechanical problem would remain.
\item \emph{Data vintages.} The geopolitical risk index is a single 2026 vintage, and bank statements
  and news come from the archive of one provider as it stands today. Items removed later cannot be
  recovered.
\item \emph{Short post-upgrade window.} The market-upgrade test has a single post-announcement window
  of about a year; its placebos vary widely, so the test has low power.
\item \emph{Informative censoring.} Swings of VN30 stocks are censored when the stock leaves the
  index, and exits can follow poor performance. Censoring affects only \CensoredPct\% of rows.
\item \emph{Pre-registration by log.} The design was fixed in a dated decision log before the test
  years were used, not in a registry with an external timestamp.
\end{itemize}

\paragraph{What the results do not show} The study makes no causal claim about the effect of any
feature on reversals. A feature that adds no forecasting skill may still matter economically if its
effect is already priced at the decision time. And the absence of trading value under our rules
does not rule out value under other rules, for example with short sales once they are allowed.
